\subsection{A Call That Changes an Operation's Validity}
\label{sec:motivating-example}
Figure~\ref{fig:motivating-example} illustrates why type-error detection
requires tracking the state established by a call. The \texttt{Message}
class declares \texttt{body: str | bytes}. Both encoders accept a string. The \texttt{encode} method of the
text encoder returns that string; the UTF-8 encoder returns bytes. The helper
\texttt{encode\_message} writes the encoder's result back to the same
message object and implicitly returns \texttt{None}.

\begin{figure}[!htbp]
\centering
\includegraphics[width=\linewidth]{motivatingExample.pdf}
\caption{A call changes the type state used by a subsequent operation.
Point 1 establishes a string body; point 2 relates the selected encoder to
the field update; point 3 has a bytes body, making concatenation with a
string fail. The lower-right panel reports the illustrated example's
outcomes for Mypy, Pyrefly, Pyright, and our approach, PyTT.}
\Description{The caller make_line raises ValueError when the message body
is not a string. Otherwise it calls encode_message with Utf8Encoder, then
passes msg.body plus a newline string to print. The helper writes the
encoder result into msg.body. TextEncoder returns a string and Utf8Encoder
returns bytes. Numbered points mark the string state before the call,
the encoder-dependent update, and the bytes state after the call. The
result panel shows no diagnostic for Mypy, Pyrefly, and Pyright and a type
error detected by PyTT.}
\label{fig:motivating-example}
\end{figure}

\paragraph{From the guard to the update.}
In \texttt{make\_line}, a non-string body triggers \texttt{ValueError}.
Every execution reaching point~1 has therefore passed the guard with a
string body. The subsequent call supplies that message and a
\texttt{Utf8Encoder} instance to the helper. At point~2, the helper's two
encoder cases connect the input choice to different updates: the text
encoder establishes a string body and the UTF-8 encoder establishes bytes.
Both writes satisfy the field's union annotation. The relevant difference
is the state that each call establishes for subsequent operations.

\paragraph{From the update to the failure.}
Instantiating the UTF-8 case at the shown call binds the helper's message
to the caller's object. Consequently, point~3 contains a bytes body.
The string fact at point~1 concerns the old contents and does not constrain
the updated field. Evaluation of \texttt{msg.body + '\textbackslash n'}
then raises \texttt{TypeError} before \texttt{print} is invoked. With
\texttt{make\_line} as the analysis entry, that exception escapes the
function. Using \texttt{TextEncoder} at the same call would instead
establish a string body and permit the concatenation.

This reasoning distinguishes three representations. The helper's return
type, \texttt{None}, gives no description of the field update. A set of
possible field types, \texttt{str | bytes}, captures the possible contents.
A relation between the selected encoder, incoming body, and outgoing body
identifies the contents established in the actual calling context.
Preserving that relation lets the analyzer invalidate the old field fact
and check the addition using its new value.

The result panel in Figure~\ref{fig:motivating-example} identifies our
approach as \emph{PyTT}. For the illustrated example, it records no
diagnostic from Mypy, Pyrefly, or Pyright and a type error detected by PyTT;
the associated tool versions and configurations are xxx. The mechanism
illustrated by this result is the propagation of the encoder-dependent
write to the operation that consumes the field. Section~\ref{sec:evaluation}
evaluates detection across the benchmark under the stated configurations.
